\section{Methodology}
\label{sec:method}

\paragraph{Throughput.} Each run skips 30 warm-up steps (compilation, allocator growth), then measures
for 5 minutes, converted to a step count that all ranks agree on. Step times come from CUDA events on
rank~0. The official throughput is the \emph{median} of the per-step global token rates, reported per
GPU (tokens/s per GPU). Spread within a run is the 10th--90th percentile of step times.

\paragraph{MFU} $= \text{FLOPs per token} \times \text{tokens/s per GPU} / \text{peak}$, with TorchTitan's
FLOP count (6 $\times$ active parameters plus attention terms, no causal halving) and a dense BF16 peak of
989\,TFLOPS per H100. MFU is derived from a measured rate and an analytic FLOP count.

\paragraph{Placement families} (at a fixed node count):
\textbf{S}, all nodes on one switch (ihpc316, the 12-node leaf without CPU nodes);
\textbf{H}, 3 nodes on each of $N/3$ switches (3-node leaves; big leaves only beyond 13 switches);
\textbf{O}, 1 node on each of $N$ switches (up to 16);
\textbf{P}, the fewest switches (largest leaves first);
\textbf{ALL}, every healthy node.
S, H and O are compared at 3, 6, 9 and 12 nodes; at scale, H15--P15, H27--P27 and H48--P48 (variant
names give the node count). Placements are written as leaves $\times$ nodes per leaf, e.g.\ $9\times3$.

\paragraph{Expert-group placement} (GLM-5, Rorqual; an expert group is an expert-parallel group of 8
GPUs = 2 nodes). Three layouts on the same 52 nodes (208 GPUs), the largest set with an even number of
nodes on every leaf ($13\times2$ on the 3-node leaves, $12+8+6$ on the big ones): \emph{aligned}, each
group's 2 nodes on one switch; \emph{crossed}, each group's 2 nodes on two switches; \emph{mixed}, about
one group in four crossing, close to the share of a production job on every node. Megatron forms groups
from consecutive ranks, so only the node order differs: model, scale, nodes and data parallelism are the
same. A fourth run uses every healthy node in Megatron's default order (\emph{mixed, every node}): the
largest achievable job, with more nodes and some groups crossing.

Node lists are computed from the healthy nodes of the allocation by \texttt{tt\_pick\_nodes.py}.

\paragraph{Traffic isolation.} Placement comparisons are only valid between runs exposed to the same
external traffic. Three levels: \emph{alone on the cluster} (the October~30 window); \emph{reserved, one
switch} (every node in our reservation and on one leaf: other users' traffic never shares our links);
\emph{shared links} (the run crosses the spine, or its nodes are not reserved: other users' traffic may
share its links). Before the window, any run spanning switches is in the last level, so preliminary
comparisons between one switch and several mix placement with external traffic; they are not compared.

\paragraph{Health checks.} Before any run, every node of the allocation is checked (GPU count, MIG, ECC,
InfiniBand link state and rate, DCGM diagnostics, NCCL bandwidth within the node, through its NIC, in
pairs and per leaf). Failing nodes are excluded for the rest of the allocation; a node is used whole or
not at all.

\paragraph{Measured, estimated and synthetic values.}
\begin{itemize}[leftmargin=*, itemsep=1pt]
  \item \emph{Measured}: filled markers, solid lines.
  \item \emph{Estimated} (extrapolated or model-based, e.g.\ compute time scaled from a 16-layer run):
        dashed lines or hatched bars, and ``est.'' in the label.
  \item \emph{Synthetic} (placeholder for an experiment not run yet): hollow markers, and the figure says so.
\end{itemize}
Estimates never replace measurements; both are kept in the data.

\paragraph{Data.} Everything is regenerated from raw results by one command
(\texttt{scripts/analysis/analyze.py all}): milabench run data (one JSONL file per run, including failed
runs), health-check outputs, and node telemetry (\texttt{benchmate.nodemon}: GPU clocks, throttle reasons,
power, temperature and duty cycle, XID and ECC errors, NVLink, InfiniBand, CPU clocks, host memory and PCIe,
sampled once a second on every node of the job). Formats: \texttt{FORMAT.md}.
